problem:  
Let \((M^n,g)\) be a complete Riemannian manifold with nonnegative Ricci curvature, and let \((N^m,h)\) be a compact manifold with strictly negative sectional curvature. Suppose \(f:M\to N\) is a finite-energy harmonic map. Classical Bochner-type results imply that if \(\operatorname{Ric}_M>0\) somewhere, then \(f\) must be constant.  

Consider instead the borderline case \(\operatorname{Ric}_M\ge 0\):  

Assume that \(f(M)\) has topological dimension \(<m\) but does **not** lie in any totally geodesic submanifold of \(N\). Must \(f\) necessarily be constant?  

Equivalently: does the absence of lower-dimensional totally geodesic submanifolds in \(N\) force that *any nonconstant harmonic map from a nonnegatively Ricci curved domain has full-dimensional image*?  

Why is it a "good" problem:  
This refined problem eliminates potential degeneracies while sharpening a genuine borderline rigidity phenomenon. Classical Bochner arguments control \(|df|^2\), but not the *rank distribution* of \(df\) or the *dimension of the image*. The question thus asks whether the analytic rigidity implied by nonnegative Ricci curvature, together with strict convexity features of negatively curved targets, extends to prohibit partial collapse of the differential.  

The problem is conceptually simple—whether “partial collapse without total triviality” can occur—but its resolution would likely require forging new links among:  

1. **Analytic curvature control** (Bochner-type inequalities for harmonic maps),  
2. **Geometric measure theoretic and stratification methods** (understanding which dimensions images can occupy), and  
3. **Global geometric rigidity** (effects of the absence of totally geodesic submanifolds).  

This interplay between Riemannian geometry, geometric analysis, and measure-theoretic geometry exemplifies a “narrow entrance, profound depth” problem. A positive solution would expose new rigidity phenomena for harmonic maps, while a counterexample would reveal unexpected flexibility in the borderline nonnegative Ricci case—both outcomes would significantly enrich our understanding of geometric rigidity.